\section{Internetworking + DNS}
\subsection{Internet Protocol (IP)}
\subsubsection*{Was ist das Internet?}
Das Internet ist ein Netzwerk von Netzwerken - also ein Zusammenschluss vieler kleinerer, oft sehr unterschiedlicher Netzwerke:
\begin{itemize}
    \item \fatsf{Internet = \enquote{Inter}-Netzwerk:} Ein globales System, das viele kleinere Netzwerke miteinander verbindet
    \item \fatsf{Router als zentrale Verbindungsstellen:} Router übernehmen das Weiterleiten (Forwarding) von Datenpaketen zwischen Netzwerken
    \item \fatsf{Kleinere Netzwerke = LANs:} Oft basieren sie auf dem IEEE-802-Standard (z.B. Ethernet, WLAN)
    \item \fatsf{Zwei Arten der Kommunikation:}
    \begin{itemize}
        \item \fatsf{Lokal (Local):} Kommunikation innerhalb eines LANs, ohne Router
        \item \fatsf{Remote (Fern):} Kommunikation über Router hinweg - über das Internet
    \end{itemize}
\end{itemize}
\subsubsection*{Was braucht man, damit das Internet funktioniert?}
Damit verschiedene Netzwerke miteinander kommunizieren können, sind mehrere Komponenten nötig:
\begin{enumerate}
    \item \fatsf{Logische Adressierung}
    \begin{itemize}
        \item Physikalische Adressen (z.B. MAC-Adressen von IEEE 802) sind nicht ausreichend
        \item \fatsf{IP-Adressen} (z.B. 192.168.1.1) dienen als logische Adressen, unabhängig von der zugrundeliegenden Hardware
    \end{itemize}
    \item \fatsf{Zuordnung von logischen zu physikalischen Adressen}
    \begin{itemize}
        \item z.B. mit \fatsf{ARP (Address Resolution Protocol)}
        \item Wird benötigt, um Pakete im lokalen Netz korrekt zuzustellen
    \end{itemize}
    \item \fatsf{Mechanismus zur Unterscheidung: lokal oder remote?}
    \begin{itemize}
        \item Beim Senden eines Pakets muss entscheiden werden:
        \begin{itemize}
            \item Ist der Empfänger im selben LAN?
            \item Oder muss das Paket über das Internet weitergeleitet werden?
        \end{itemize}
        \item Diese Entscheidung trifft meist der Router anhand von Routing-Tabellen
    \end{itemize}
    \item \fatsf{Unterstützung für viele Technologien}
    \begin{itemize}
        \item Das Internet muss verschiedenste Übertragungstechniken unterstützen:
        \begin{itemize}
            \item Ethernet, WLAN, Glasfaser, Mobilfunk usw.
        \end{itemize}
        \item Dazu gehören auch Funktionen wie:
        \begin{itemize}
            \item \fatsf{Fragmentierung} (Aufteilen großer Pakete)
            \item \fatsf{Diagnose} (z.B. ICMP für Fehlermeldungen)
        \end{itemize}
        $\rightarrow$ All das wird durch das Internetprotokoll \fatsf{IP (Internet Protocol)} möglich gemacht
    \end{itemize}
\end{enumerate}
\subsubsection*{IPv4 - Das klassische Internetprotokoll}
\begin{itemize}
    \item \fatsf{Historie}
    \begin{itemize}
        \item Ursprünglich definiert in RFC 760 (Januar 1980)
        \item Weltweit im Einsatz - lange das Rückgrat des Internets
        \item Verwendet 32-Bit-Adressen $\rightarrow$ max. ca. 4,3 Mrd. eindeutige Adressen
    \end{itemize}
    \item \fatsf{IPv4 Adressierung}
    \begin{itemize}
        \item Beispieladresse: 192.168.1.53
        \item Darstellung in vier Dezimalblöcken (z.B. 192.168.0.1)
    \end{itemize}
    \item \fatsf{Helferprotokolle für IPv4}
    \begin{table}[H]
        \centering
        \rowcolors{2}{\thepagecolor}{fgcolor!10!\thepagecolor}
        \renewcommand{\arraystretch}{1.2}
        \begin{tabular}{ll}
            \toprule
            \fatsf{Funktion} & \fatsf{Protokoll}\\
            \midrule
            Adressvergabe & DHCP\\
            Auflösung IP $\rightarrow$ MAC & ARP (Request/Reply)\\
            Fehlermeldungen, Diagnose & ICMPv4\\
            Multicast-Gruppen & IGMP\\
            \bottomrule 
        \end{tabular}
    \end{table}
    \item \fatsf{Subnetting}
    \begin{itemize}
        \item Aufteilung eines Netzwerks in Subnetze mithilfe von Subnetzmasken, z.B. 255.255.255.0
        \item Dient zur Unterscheidung von Netzanteil in der Adresse
    \end{itemize}
    \item \fatsf{Beispielkommunikation (vereinfacht):}
    \begin{itemize}
        \item Host A will Daten an 192.168.1.42 senden
        \item Prüft: Ist Ziel im eigenen Netz?
        \item $\rightarrow$ Wenn ja: ARP-Request $\rightarrow$ MAC-Adresse von 192.168.1.42 wird gefragt
        \item $\rightarrow$ Dann: Direkt senden. Falls nein: Routing über Standard-Gateway
    \end{itemize}
\end{itemize}
\subsubsection*{Subnetzmaske}
\begin{itemize}
    \item \fatsf{Was ist eine Subnetzmaske?}\\
    Eine \fatsf{Subnetzmaske} (engl. subnet mask) trennt bei einer IP-Adresse den \fatsf{Netzwerkanteil} vom \fatsf{Hostanteil}. Sie wird gemeinsam mit einer IP-Adresse verwendet, um zu bestimmen:
    \begin{itemize}
        \item \fatsf{Zu welchem Netzwerk} eine IP-Adresse gehört
        \item \fatsf{Welche Adressen} im selben Subnetz liegen
    \end{itemize}
    \item \fatsf{Beispiel einer IPv4-Adresse und ihrere Subnetzmaske:}
    \begin{itemize}
        \item IP-Adresse: 192.168.1.10
        \item Subnetzmaske: 255.255.255.0
    \end{itemize}
    In binärer Form:
\begin{lstlisting}
IP-Adresse:     11000000.10101000.00000001.00001010
Subnetzmaske:   11111111.11111111.11111111.00000000
\end{lstlisting}   
    Die 1-Bits in der Subnetzmaske markieren den \fatsf{Netzwerkanteil}, die 0-Bits den \fatsf{Hostanteil}\\
    \fatsf{In diesem Fall:}
    \begin{itemize}
        \item Netzwerkanteil: 192.168.1
        \item Hostanteil: 10
    \end{itemize}
    Das bedeutet: Alle Geräte mit IP-Adressen von 192.168.1.1 bis 192.168.1.254 liegen im \fatsf{gleichen Subnetz}
    \item \fatsf{Wozu braucht man Subnetzmasken?}
    \begin{enumerate}
        \item \fatsf{Netzwerke unterteilen (Subnetting)}\\
        Große Netzwerke werden in kleinere Subnetze aufgeteilt, um:
        \begin{itemize}
            \item Broadcast-Verkehr zu reduzieren
            \item Sicherheit zu erhöhen
            \item IP-Adressen effizienter zu verwalten
        \end{itemize}
        \item \fatsf{Routing}\\
        Router verwenden Subnetzmasken, um zu entscheiden, ob ein Zielgerät \fatsf{im eigenen Netzwerk} liegt oder ob das Paket \fatsf{weitergeleitet} werden muss
    \end{enumerate}
    \item \fatsf{CIDR-Notation}\\
    Subnetzmasken werden oft in \fatsf{CIDR-Notation} geschrieben, z.B.: 
    \begin{grayInfoBox}
        192.168.1.10/24
    \end{grayInfoBox}
    Die Zahl 24 steht für die Anzahl der 1-Bits in der Subnetzmaske:
    \begin{grayInfoBox}
        /24 = 255.255.255.0
    \end{grayInfoBox}
    Einige gängige CIDR-Werte:
    \begin{table}[H]
        \centering
        \rowcolors{2}{\thepagecolor}{fgcolor!10!\thepagecolor}
        \renewcommand{\arraystretch}{1.2}
        \begin{tabular}{lll}
            \toprule
            \fatsf{CIDR} & \fatsf{Subnetzmaske} & \fatsf{Anzahl Hosts}\\
            \midrule
            /8 & 255.0.0.0 & 16.777.214 Hosts\\
            /16 & 255.255.0.0 & 65.534 Hosts\\
            /24 & 255.255.255.0 & 254 Hosts\\
            /30 & 255.255.255.252 & 2 Hosts\\
            \bottomrule 
        \end{tabular}
    \end{table}
    (immer abzüglich Netz- und Broadcast-Adresse)
\end{itemize}


\subsubsection*{IPv6 - Die moderne Alternative}
\begin{itemize}
    \item \fatsf{Historie \& Motivation}
    \begin{itemize}
        \item Entwickelt ab 1995, RFC 2460 erschien 1998
        \item Ziel: Rückkehr zur Einfachheit \& Autonomie der Anfangsjahre
        \item Hauptgründe für Einführung:
        \begin{itemize}
            \item End-to-End Prinzip wiederherstellen
            \item Riesiger Adressraum (128 Bit $\rightarrow$ 3,4 $\times$ $10^{38}$ Adressen)
        \end{itemize}
    \end{itemize}
    \item \fatsf{IPv6 Adressierung}
    \begin{itemize}
        \item Beispieladresse: 2001:0db8::1
        \item Hexadezimal dargestellt
        \item Automatische Konfiguration möglich (SLAAC)
        \item Trennung in Netz- und Hostanteil meist bei /64
    \end{itemize}
    \item \fatsf{Helferprotokolle für IPv6}
    \begin{table}[H]
        \centering
        \rowcolors{2}{\thepagecolor}{fgcolor!10!\thepagecolor}
        \renewcommand{\arraystretch}{1.2}
        \begin{tabular}{ll}
            \toprule
            \fatsf{Funktion} & \fatsf{Protokoll}\\
            \midrule
            Adressvergabe & DHCPv6, SLAAC\\
            Auflösung IP $\rightarrow$ MAC & NDP (NS/NA)\\
            Routinginformation & RA, RS\\
            Multicast-Management & MLD\\
            Fehlermeldungen, Diagnose & ICMPv6\\
            \bottomrule 
        \end{tabular}
    \end{table}
\end{itemize}
\subsubsection*{Warum ist der Umstieg so schwierig?}
\begin{itemize}
    \item \fatsf{Vorteile von IPv6:}
    \begin{itemize}
        \item Fast unbegrenzter Adressraum
        \item Keine Konflikte durch private IPv4-Adressen \& NAT
        \item Entlastung bei IPSs durch Wegfall komplizierter NAT-Strukturen
    \end{itemize}
    \item \fatsf{Nachteile bzw. Herausforderungen:}
    \begin{itemize}
        \item Nicht kompatibel mit IPv4 $\rightarrow$ beide Protokolle müssen parallel betrieben werden (Dual Stack)
        \item Anpassung von Netzwerk- und Sicherheitsinfrastruktur nötig
        \item Geringer Nutzen für viele Unternehmen $\rightarrow$ hohes Aufwand-Nutzen-Verhältnis
    \end{itemize}
\end{itemize}
\clearpage
\subsection{Domain Name Service (DNS)}
\subsubsection*{Was ist das Problem, das DNS löst?}
\begin{itemize}
    \item \fatsf{Ausgangspunkt}\\
    Wir haben gelernt, dass im Internet Geräte durch IP-Adressen eindeutig adressiert werden - z.B. 130.83.47.181. Aber für Menschen ist es schwer, sich solche Adressen zu merken
    $\rightarrow$ Problem: Wie kann Alice Bob kontaktieren, wenn sie seine IP-Adresse nicht kennt?
    \item \fatsf{Lösung}\\
    Das Domain Name System (DNS) löst dieses Problem:
    \begin{itemize}
        \item DNS ist ein verteilter Dienst, der Hostnamen (wie \url{www.tu-darmstadt.de}) in IP-Adressen übersetzt
        \item Man nennt das auch \enquote{Namensauflösung}
    \end{itemize}
\end{itemize}
\begin{figure}[h]
    \centering
    \includegraphics[width=0.5\linewidth]{Chapter_3_Images/DNS.jpeg}
\end{figure}
\subsubsection*{Warum braucht man verschiedene Arten von Adressen?}
Im Internet gibt es auf verschiedenen Schichten unterschiedliche Adresstypen:
\begin{table}[H]
    \centering
    \rowcolors{2}{\thepagecolor}{fgcolor!10!\thepagecolor}
    \renewcommand{\arraystretch}{1.2}
    \begin{tabular}{llll}
        \toprule
        \fatsf{Ebene} & \fatsf{Beispiel} & \fatsf{Art der Adresse} & \fatsf{Eigenschaften}\\
        \midrule
        Application Layer & \url{www.tu-darmstadt.de} & Hostname & Menschlich lesbar, hierarchisch\\
        Internet Layer & 130.83.47.181 & IP-Adresse (IPv4/IPv6) & Maschinenlesbar, hierarchisch\\
        Data Link Layer & 00:97:F1:12:34:56 & MAC-Adresse & Maschinenlesbar, flach (nicht hierarchisch)\\
        \bottomrule 
    \end{tabular}
\end{table}
$\rightarrow$ Deshalb braucht man Mechanismen, um zwischen diesen Adresstypen zu übersetzen (DNS für Hostname $\rightarrow$ IP, ARP für IP $\rightarrow$ MAC)\\
\fatsf{Hierarchical vs. Flat}
\begin{itemize}
    \item \fatsf{Hierarchisch}
    \begin{itemize}
        \item Die Adressen sind \fatsf{in Ebenen oder Schichten organisiert}
        \item Sie können in Teile zerlegt werden, die jeweils eine Bedeutung haben (z.B. bei einer IP-Adresse: Netzwerkanteil und Hostanteil)
        \item Das hilft, die Adressen effizient zu strukturieren, zu verwalten und zu routen
        \item Beispiel: IP-Adressen - der Teil links bestimmt das Netzwerk, der Teil rechts das Gerät im Netzwerk
    \end{itemize}
    \item \fatsf{Flach:}
    \begin{itemize}
        \item Die Adressen sind \fatsf{eindimensional ohne Struktur oder Unterteilung}
        \item Jedes Gerät hat eine eindeutige, aber nicht hierarchisch gegliederte Adresse
        \item Es gibt keine \enquote{Untergruppen} oder logische Schichten in der Adresse
        \item Beispiel: MAC-Adressen - jede Adresse ist einzigartig, aber es gibt keine Teile, die eine Struktur oder Organisation anzeigen
    \end{itemize}
\end{itemize}
\subsubsection*{High Order and Low Order}
\begin{enumerate}
    \item \fatsf{Grundidee:}
    \begin{itemize}
        \item \fatsf{High Order} bezeichnet den \fatsf{höherwertigen, übergeordneten} Teil eines Identifikators oder einer Adresse\\
        Er gibt eine \fatsf{breite, grobe oder übergeordnete Zuordnung} an (z.B. ein Netzwerk, eine Organisation oder ein Hersteller)
        \item \fatsf{Low Order} bezeichnet den \fatsf{niederwertigen, spezifischeren} Teil eines Identifikators. Er steht für \fatsf{feinere, detailliertere oder individuelle Zuordnung} (z.B. ein einzelner Host, ein Interface oder ein Gerät innerhalb der übergeordneten Struktur)
    \end{itemize}
    \item \fatsf{Beispiel:}
    \begin{table}[H]
        \centering
        \rowcolors{2}{\thepagecolor}{fgcolor!10!\thepagecolor}
        \renewcommand{\arraystretch}{1.2}
        \begin{tabular}{lp{6cm}p{6cm}}
            \toprule
            \fatsf{Identifikator} & \fatsf{High Order (übergeordnet)} & \fatsf{Low Order (fein)}\\
            \midrule
            \fatsf{Hostname} & Domain, die von einer Registrierungsstelle vergeben wird, z.B. .com, .edu oder tu-darmstadt.de (z.B. bei seemoo.tu-darmstadt.de: .de und tu-darmstadt sind High Order) & Lokaler Hostname, vom Administrator vergeben, z.B. \texttt{www}, \texttt{cms-sip02} (\enquote{www} ist Low Order, da spezifischer Host in der Domain)\\
            \fatsf{IP-Adresse} & Netzwerkteil (Präfix), vergeben von ICANN, Regionalen Internetregistern (RIR) oder ISPs (z.B. bei 130.83.47.181 ist 130.83 das Netzwerk) & Host-Teil der IP-Adresse, einem Gerät im Netzwerk zugewiesen (z.B. per DHCP oder manuell) (z.B. 47.181 ist Low Order)\\
            \fatsf{MAC-Adresse} & Herstellerkennung (Organizationally Unique Identifier, OUI), 24 Bits von IEEE an Hersteller vergeben (z.B. 00-15-C5 gehört z.B. zu IBM, Apple) & Gerätespezifische Schinittstellenkennung, 24 Bits, vom Hersteller vergeben (z.B. 49-04-A9 ist Low Order für die Netzwerkkarte)\\
            \bottomrule 
        \end{tabular}
    \end{table}
\end{enumerate}
\subsubsection*{ICANN (Internet Corporation for Assigned Names and Numbers)}
\begin{itemize}
    \item \fatsf{ICANN} ist eine \fatsf{nicht-kommerzielle Organisation}, die weltweit zentrale Aufgaben für das Funktionieren des Internets übernimmt
    \item \fatsf{Hauptaufgaben von ICANN:}
    \begin{enumerate}
        \item \fatsf{Domain Name System (DNS) verwalten:}\\
        ICANN koordiniert die Vergabe von Domainnamen (z.B. .com, .org, .de) und überwacht die sogenannten \fatsf{Top-Level-Domains (TLDs)}
        \item \fatsf{Zuteilung von IP-Adressen:}\\
        Sie delegiert IP-Adressblöcke anregionale Organisationen (wie RIPE für Europa), die diese weiter an Provider und Unternehmen vergeben
        \item \fatsf{Verwaltung von Root-Servern:}\\
        ICANN verwaletet zusammen mit anderen Organisationen das Root-Server-System des Internets, das die Grundlage für DNS-Anfragen ist
        \item  \fatsf{Zuständigkeit für Protokollparameter:}\\
        In Zusammenarbeit mit der \fatsf{IETF (Internet Engineering Task Force)} sorgt ICANN für die Koordination technischer Standards und Protokollnummern
    \end{enumerate}
\end{itemize}
\subsubsection*{RIR (Regional Internet Registry)}
\begin{itemize}
    \item Ein RIR ist eine Organisation, die im Auftrag von \fatsf{ICANN} IP-Adressen und autonome Systemnummern (\fatsf{ASNs}) in einer bestimmten geographischen Region verwaltet und zuteilt
    \item \fatsf{Die 5 RIRs weltweit:}
    \begin{table}[H]
        \centering
        \rowcolors{2}{\thepagecolor}{fgcolor!10!\thepagecolor}
        \renewcommand{\arraystretch}{1.2}
        \begin{tabular}{ll}
            \toprule
            \fatsf{Region} & \fatsf{RIR-Name}\\
            \midrule
            Europa, Naher Osten, Zentralasien & RIPE NCC\\
            Nordamerika & ARIN\\
            Lateinamerika \& Karibik & LACNIC\\
            Afrika & AFRINIC\\
            Asien-Pazifik & APNIC\\
            \bottomrule 
        \end{tabular}
    \end{table}
    \item \fatsf{Hauptaufgaben eines RIR:}
    \begin{enumerate}
        \item \fatsf{Zuteilung von IP-Adressen} (IPv4 \& IPv6) und \fatsf{AS-Nummern} an Internetdienstanbieter (ISPs), Unternehmne, Universitäten usw.
        \item \fatsf{Führen öffentlicher Datenbanken}, um festzuhalten, wem welche IP-Adressen gehören (z.B. für Whois-Abfragen)
        \item \fatsf{Förderung technischer Standards}, Schulungen, Sicherheit und verantwortungsvoll IP-Nutzung
        \item \fatsf{Koordination mit ICANN} und anderen RIRs, um globale Adresspolitik umzusetzen
    \end{enumerate}
    \item \fatsf{Beispiel:}\\
    Wenn z.B. ein Internetprovider in Deutschland neue IP-Adressen braucht, beantragt er sie bei \fatsf{RIPE NCC}
\end{itemize}
\subsubsection*{Implementing DNS: Alte Methode (Hosts-Datei)}
\begin{itemize}
    \item \fatsf{Was:}\\
    Eine lokale Datei (\texttt{/etc/hosts} auf Unix, \texttt{$C:\textbackslash \cdots \textbackslash hosts$} auf Windows) enthält Zuordnungen von Hostnamen zu IP-Adressen
    \item \fatsf{Wie:}\\
    Bevor DNS befragt wird, prüft das System diese Datei - findet es einen Eintrag, wird die IP direkt verwendet
    \item \fatsf{Früher:}\\
    War notwendig, als das Internet noch klein war. Alle Hosts mussten ihre Datei manuell aktuell halten
    \item \fatsf{Warum heute ungeeignet:}
    \begin{itemize}
        \item Nicht skalierbar (zu viele Hosts)
        \item Keine automatische Aktualisierung
        \item Sicherheitsrisiken (manipulierbar)
        \item Fehleranfällig
    \end{itemize}
    \item \fatsf{Heute genutzt für:}
    \begin{itemize}
        \item Lokale Tests
        \item Webseiten blockieren
        \item Entwicklung ohne DNS
    \end{itemize}
\end{itemize}
\subsubsection*{Warum kein zentrales DNS?}
\begin{itemize}
    \item DNS ist absichtlich dezentral aufgebaut:
    \begin{itemize}
        \item Skalierbarkeit: Ein zentrales System wäre überlastet
        \item Ausfallsicherheit: Single Point of Failure (SPOF) wird vermieden
        \item Verwaltungshoheit: Jede Organisation kann eigene DNS-Einträge verwalten
    \end{itemize}
\end{itemize}
\subsubsection*{Welche Arten von DNS-Einträge (Resource Records) gibt es?}
\begin{table}[H]
    \centering
    \rowcolors{2}{\thepagecolor}{fgcolor!10!\thepagecolor}
    \renewcommand{\arraystretch}{1.2}
    \begin{tabular}{ll}
        \toprule
        \fatsf{Typ} & \fatsf{Beschreibung}\\
        \midrule
        A & IPv4-Adresse\\
        AAAA & IPv6-Adresse\\
        CNAME & Alias auf anderen Hostnamen\\
        MX & Mailserver (Mail eXchanger)\\
        NS & Nameserver für eine Domain\\
        TXT & Freitext - heute z.B. für SPF, DKIM etc.\\
        SOA & Informationen über Zone-Verwaltung\\
        \bottomrule 
    \end{tabular}
\end{table}
\subsubsection*{Aufbau eines DNS Resource Records (RR):}
\begin{figure}[h]
    \centering
    \includegraphics[width=0.9\linewidth]{Chapter_3_Images/DNS on the Wire.jpeg}
\end{figure}
\begin{itemize}
    \item \fatsf{TTL:} Time To Live in Sekunden, wie lange der RR im Cache bleibt
    \item \fatsf{Class:} Klasse des RR, meist IN für Internet
    \item \fatsf{RL: RD LENGTH:} Länge der darauffolgenden RDATA in Bytes
    \item \fatsf{RDATA:} Daten des RR, z.B. IP-Adresse für A, Mailserver für MX
\end{itemize}
\subsubsection*{Wie funktioniert eine DNS-Anfrage?}
\begin{enumerate}
    \item Du gibst eine Domain im Browser ein
    \item Der \fatsf{Stub Resolver} deines Betriebssystems schickt die Anfrage an den \fatsf{rekursiven Resolver}
    \item Der rekursive Resolver sucht die IP-Adresse, indem er verschiedene DNS-Server abfragt
    \item Er holt sich die Antwort vom \fatsf{autoritativen DNS-Server} der jeweiligen Domain
    \item Der rekursive Resolver speichert die Antwort im Cache für spätere Anfragen
    \item Die IP-Adresse wird an den Stub Resolver zurückgegeben
    \item Dein Computer kann die Webseite jetzt mit der IP-Adresse erreichen
\end{enumerate}
\begin{figure}[h]
    \centering
    \includegraphics[width=0.7\linewidth]{Chapter_3_Images/DNS query example.jpeg}
\end{figure}
\subsubsection*{DNS und Sicherheit}
\begin{itemize}
    \item \fatsf{DNS ist essentiell} für das Internet, weil es Domain-Namen in IP-Adressen übersetzt. Ohne DNS funktionieren viele Dienste nicht
    \item Ein Ausfall von DNS (z.B. durch falsches Routing wie bei Facebook 2021) kann ganze Webseiten lahmlegen
    \item DNS betrifft wichtige Sicherheitsaspekte:
    \begin{itemize}
        \item \fatsf{Vertraulichkeit:} DNS-Anfragen verraten, welche Seiten Nutzer besuchen
        \item \fatsf{Integrität:} Manipulierte DNS-Daten können Nutzer zu falschen Seiten umleiten
        \item \fatsf{Verfügbarkeit:} Wenn DNS nicht funktioniert oder blockiert wird, sind Webseiten nicht erreichbar
    \end{itemize}
    \item DNS ist oft ungeschützt und anfällig für Überwachung, Manipulation und Zensur
\end{itemize}